\section{Métricas del agente y controles de calidad}
\subsection{Métricas clave}
Además de la tasa de éxito en SWE-bench, es necesario prestar atención a otras métricas: la tasa de intervención humana, el tiempo de espera promedio, la cobertura de pruebas y la tasa de diferencias. Estas métricas ayudan al equipo a determinar si el agente realmente reduce la carga de ingeniería en lugar de agregar costos adicionales.

\begin{importantbox}{Desglose de métricas}
Graham recomienda prestar atención a cuatro dimensiones: 1) tasa de éxito (aprobado/reprobado), 2) tasa de intervención humana (cuántas veces se requiere pausar), 3) tiempo de ejecución (si está dentro de un rango aceptable), 4) interpretabilidad (si la salida ofrece una explicación razonable).
\end{importantbox}

\subsection{Nodos de control de calidad}
En el pipeline del agente, los controles de calidad pueden colocarse en los nodos Plan, Act o Verify: Plan requiere una revisión humana de los pasos generados, la etapa Act necesita una revisión de los comandos por ejecutar y la etapa Verify consiste en pruebas y análisis estático. Solo cuando se superan todos los controles se permite enviar el patch.

\begin{table}[H]
\centering
\begin{tabular}{p{0.3\textwidth} p{0.33\textwidth} p{0.33\textwidth}}
\toprule
Nodo de control & Criterio & Respuesta \\
\midrule
Plan & Si los pasos del plan son válidos y si implican operaciones con permisos & Revisión humana o adición de un prompt \\
\midrule
Act & Si el comando involucra recursos sensibles (bases de datos, servicios en la nube) & Pausar y solicitar aprobación \\
\midrule
Verify & Si las pruebas o el análisis estático pasan en su totalidad & Revertir, agregar pruebas o someter a revisión humana \\
\bottomrule
\end{tabular}
\caption{Disposición típica de los controles de calidad en el flujo del agente.}
\end{table}

\begin{warningbox}{No te fijes solo en la tasa de éxito}
Una tasa de éxito alta puede significar que el agente está limitado a tareas de bajo riesgo; en un entorno de producción real también es necesario considerar de manera conjunta la evaluación humana y el costo de ejecución.
\end{warningbox}

\subsection{Resumen de la sección}
\begin{itemize}
    \item La tasa de éxito es solo una de las métricas; hay que evaluarla junto con la tasa de intervención, el tiempo de ejecución y otros factores.
    \item Los controles de calidad ayudan a intervenir en los nodos clave, evitando que el agente publique directamente e introduzca riesgos.
\end{itemize}
